% paper/sections/12g_time_integration.tex
%
% Tier VI: 制約・否定テスト
%   U9 = DCCD-on-pressure 禁止則 否定テスト
% Tier VII: coupling primitive (single mechanism)
%   U7 = BF static droplet 1-step (Laplace 圧力誤差)
%        (a) 円形界面静止 Δp = σκ vs 計算値 (BF 演算子整合性)
%        (b) ρ/μ 補間 (arithmetic/harmonic/VOF) 面値精度
% ═══════════════════════════════════════════════════════════════════════════════
\subsection{Tier VI：制約・否定テスト}
\label{sec:verify_constraint}

Part 2 の設計選択には，「ある操作をしてはいけない」という
\textbf{制約条件}が含まれる．本 Tier では，特に重要な
\textbf{DCCD-on-pressure 禁止則}（§~\ref{sec:dccd_pressure_nofilt}）を
\textbf{否定テスト}（違反すると静止液滴に寄生流れが噴出する）として検証する．

% ─────────────────────────────────────────────────────────────────────────────
\subsubsection{U9：DCCD-on-pressure 禁止則 否定テスト}
\label{sec:U9_dccd_pressure_prohibition}
% ─────────────────────────────────────────────────────────────────────────────

\noindent\textbf{目的：}
§~\ref{sec:dccd_pressure_nofilt} で導入された「圧力場に DCCD フィルタを適用してはならない」
という設計制約を，違反した場合の不安定性増大を観測することで\textbf{否定的に}検証する．

\noindent\textbf{Part 2 対応 primitive：} \#20 (DCCD-pressure 禁止則)．

\noindent\textbf{下流連関：}
Static droplet 200-step stability（§~\ref{sec:twophase_static_longterm}）．

\noindent\textbf{手順とパラメタ：}
\begin{itemize}[noitemsep]
  \item 静止液滴（$N = 64$，$R = 0.25$，中心 $(0.5, 0.5)$，
        $\rho_l/\rho_g = 1000$，$\sigma = 1$）
  \item 2 ケース比較：
        (i) DCCD filter OFF（規則遵守，本論文の標準設定），
        (ii) DCCD filter ON（$\varepsilon_d = 0.05, 0.25$，禁止則違反）
  \item $\Delta t = 10^{-3}$，200 ステップで $\max|\bu|$ の時間履歴を計測
\end{itemize}

\noindent\textbf{結果（要約）：}
製造解 $p = \sin(2\pi x)\sin(2\pi y)$（周期境界）に対する診断試験では，
フィルタなしの CCD ラプラシアンは厳密解に対して $\Ord{h^{6.0}}$ で収束する一方，
DCCD フィルタ（$\varepsilon_d = 0.25$）を適用すると
$\nabla^2 p - \nabla^2 \tilde{p}$（$\tilde{p}$: フィルタ後）の $L_\infty$ ノルムが
$\Ord{h^{2.0}}$ の発散誤差を生じる（$N = 32$: $1.51$，$N = 256$: $2.38 \times 10^{-2}$）．
$\varepsilon_d = 0.05$ でも同じ $\Ord{h^{2.0}}$ スケーリング．
この誤差は速度補正で
$\nabla \cdot \bu^{n+1}_\mathrm{filtered} = (\nabla^2 p - \nabla^2 \tilde{p})/\Delta t \neq 0$
となり非圧縮条件を $\Ord{h^2}$ で破壊する．
静止液滴 200-step 試験では filter ON ケースで寄生流れが急速に増大し，
filter OFF ケースは Laplace 圧周辺の機械精度近傍誤差を維持する．

\noindent\textbf{結論：}
DCCD-on-pressure 禁止則の必要性を $\Ord{h^2}$ 発散誤差として定量実証．\textbf{合格}（否定テスト成立）．

% ═══════════════════════════════════════════════════════════════════════════════
\subsection{Tier VII：coupling primitive（single mechanism）}
\label{sec:verify_coupling}

Tier VII は「複数 primitive の最小結合」を 1 ステップで検証する Tier である．
本 Tier の単体テスト U7 は BF（Balanced Force）演算子整合性を
静止液滴 1-step（Laplace 圧誤差）で測る．
これは Part 2 全体で\textbf{最重要}の coupling primitive である：
寄生流れ（parasitic current）の発生は BF 演算子不整合の直接的指標であり，
第~\ref{sec:verification}章 の Static droplet（§~\ref{sec:twophase_static_longterm}）
全件の前提となる．

% ─────────────────────────────────────────────────────────────────────────────
\subsubsection{U7：BF static droplet 1-step (Laplace 圧力誤差)}
\label{sec:U7_bf_static_droplet}
% ─────────────────────────────────────────────────────────────────────────────

\noindent\textbf{目的：}
§~\ref{sec:balanced_force} の BF（Balanced Force）整合性条件——
圧力勾配 $\nabla p^{n+1}$ と表面張力 $\sigma\kappa\nabla\psi$ に同一の
離散微分演算子を用いる——が，1-step Laplace 圧誤差として
寄生流れを抑制することを検証する．

\noindent\textbf{Part 2 対応 primitive：}
\#18 (BF discrete operator matching)，\#12 ($\rho/\mu$ 補間規則)．

\noindent\textbf{下流連関：}
Static droplet (Laplace, parasitic) §~\ref{sec:twophase_static_longterm}．

% ───────────────
\paragraph{(a) 円形界面静止 $\Delta p = \sigma\kappa$ vs 計算値（BF 演算子整合性）}
\label{sec:U7_laplace_pressure}

円形液滴（$R = 0.25$，中心 $(0.5, 0.5)$，$\sigma = 1$，重力なし）に対して
NS ソルバを 1 ステップ実行し，PPE で計算された圧力場 $p^1$ から
内外圧力差 $\Delta p_\mathrm{num}$ を抽出する．
解析解 $\Delta p_\mathrm{exact} = \sigma\kappa = \sigma/R = 4.0$ との比較．

\noindent\textbf{2 演算子整合性のテスト：}
\begin{enumerate}[noitemsep]
  \item \textbf{整合} (Match)：$\nabla p$ と $\nabla\psi$ にともに CCD $\Ord{h^6}$
        （または FCCD face-average $\mathcal{G}^\text{adj}$ 非一様格子）
  \item \textbf{不整合} (Mismatch)：$\nabla p$ に CCD，$\nabla\psi$ に 2 次 FD
\end{enumerate}
$\rho_l/\rho_g = 1000$, $\mathit{We} = 1$，$N = 32, 64, 128$．

\noindent\textbf{結果：}
整合ケースでは $\Delta p_\mathrm{num} \approx 4.0$ を全格子で相対誤差
$0.16$--$0.31\%$ で達成，$\max|\bu^1|$ は機械精度近傍に留まる．
不整合ケースでは $\max|\bu^1|$ が整合ケース比 $10^3$--$10^6$ 倍に増幅
（演算子不一致による寄生流れ；§~\ref{sec:balanced_force} の理論予測と一致）．

% ───────────────
\paragraph{(b) $\rho/\mu$ 補間（arithmetic/harmonic/VOF）面値精度}
\label{sec:U7_rho_mu_interp}

セル中心物性場 $\rho, \mu$ から face $i+1/2$ への補間規則
（arithmetic mean,harmonic mean,VOF（volume-of-fluid weighted））を
1D 平面界面で比較し，face 値の $L_\infty$ 誤差と Laplace 圧計算への影響を測定する．
全 two-phase 第~\ref{sec:verification}章 ベンチマークで使用される primitive．

\noindent\textbf{結果：}
arithmetic と VOF は界面で $\Ord{h^1}$（界面位置の解像度依存），
harmonic は $1/\rho$ 平均化により高密度比で安定．
高 $\rho_l/\rho_g$ 設定では harmonic が推奨される（§~\ref{sec:rho_mu_interpolation} 参照）．

\noindent\textbf{結論：}
BF 整合演算子で Laplace 圧再現相対誤差 $0.16$--$0.31\%$；
不整合では寄生流れが $10^3$--$10^6$ 倍に増幅．
$\rho/\mu$ 補間規則は密度比に応じて harmonic を選択することで安定．\textbf{合格}．

\medskip
以上で全 Tier I--VII の単体検証（U1--U9）が完了した．
次節（§~\ref{sec:verify_summary}）に 9 テストの pass/fail 総括と
第~\ref{sec:verification}章 への接続を示す．
